%%% Local Variables:
%%% mode: LaTeX
%%% TeX-master: "../main"
%%% End:

\section{Task C: Vary the Duty Cycle}

\subsection{a) Calculations:}

Gegeben ist eine feste Zykluszeit von \(T = 2000\,\mathrm{ms}\) und drei Duty Cycles.

\begin{table}[h]
	\centering
	\begin{tabular}{|l|c|c|c|c|}
		\hline
		             & \textbf{Step 1} & \textbf{Step 2} & \textbf{Step 3} & \textbf{Generic}               \\ \hline
		Cycle Time   & 2000 ms         & 2000 ms         & 2000 ms         & $t_{Cycle}$                    \\ \hline
		Duty Cycle   & 70\%            & 50\%            & 30\%            & $D$                            \\ \hline
		LED ON Time  & 1400 ms         & 1000 ms         & 600 ms          & $t_{on} = D \cdot t_{Cycle}$   \\ \hline
		LED OFF Time & 600 ms          & 1000 ms         & 1400 ms         & $t_{off} = t_{Cycle} - t_{on}$ \\ \hline
	\end{tabular}
\end{table}

\subsection{b) Implementation:}

Der Duty Cycle wird in drei Stufen variiert, während die Zykluszeit konstant bleibt. Für jede Stufe werden Ein- und Ausschaltzeit berechnet und die LED 15-mal geblinkt.

\begin{lstlisting}
  /* USER CODE BEGIN 2 */
  const float CycleTime = 2000;
  const float dutySteps[3] = {0.7, 0.5, 0.3};
  /* USER CODE END 2 */

  uint32_t OnTime;
  uint32_t OffTime;

  /* Infinite loop */
  /* USER CODE BEGIN WHILE */
  while (1)
  {
    /* USER CODE END WHILE */
    for(int i = 0; i < 3; i++)
    {
      float duty = dutySteps[i];

      OnTime = (uint32_t)(CycleTime * duty);
      OffTime = (uint32_t)(CycleTime - OnTime);

      // BlinkLED
      for(int i=0;i<15;i++)
      {
        GPIOJ->ODR |= SEGDP_Pin;
        HAL_Delay(OnTime);
        GPIOJ->ODR &= ~SEGDP_Pin;
        HAL_Delay(OffTime);
      }
    }
    /* USER CODE BEGIN 3 */
  }
  /* USER CODE END 3 */
\end{lstlisting}

\subsection{c) Discussion:}

Die Struktur mit einem Array für die Duty Cycles ist übersichtlich und leicht erweiterbar. Die Trennung von Berechnung und Ausführung vereinfacht die Wiederverwendung.
